% Estilo comum das figuras de vetores (capítulo 8).
\usepackage{tikz}
\usetikzlibrary{arrows.meta, positioning, calc, decorations.pathreplacing}
\definecolor{navy}{RGB}{23,45,80}
\definecolor{laranja}{RGB}{250,110,29}
\definecolor{cinzalogo}{RGB}{107,114,126}
\definecolor{azul}{RGB}{205,222,247}
\tikzset{
  casa/.style={draw=navy, line width=0.9pt, minimum width=8mm, minimum height=7mm,
               inner sep=0pt, font=\sffamily\small},
  pronta/.style={casa, fill=azul},
  parte/.style={casa, fill=laranja!12},
  marcada/.style={casa, fill=laranja!45, draw=laranja},
  fora/.style={casa, draw=cinzalogo!60, dashed, fill=white, text=cinzalogo!60},
  indice/.style={font=\sffamily\scriptsize, text=cinzalogo},
  explica/.style={font=\sffamily\small, text=navy, anchor=east},
  legenda/.style={font=\sffamily\scriptsize, text=cinzalogo},
  rotulo/.style={font=\sffamily\small\bfseries, text=navy},
  seta/.style={-{Stealth[length=2.2mm]}, navy, line width=0.9pt},
  chave/.style={decorate, decoration={brace, amplitude=4pt, mirror}, navy, line width=0.8pt},
}
% \vetor{nome}{x}{y}{valor/estilo, valor/estilo, ...}: casas lado a lado, chamadas nome0, nome1...
\newcommand{\vetor}[4]{%
  \foreach \valor/\estilo [count=\i from 0] in {#4} {
    \node[\estilo] (#1\i) at ($(#2, #3) + (\i*0.8, 0)$) {\valor};
  }
}
